%% GenABR: A Tiered, Spatially-Aware Adaptive Bitrate Streaming System
%% IEEE Conference Format — IEEEtran
%%
%% This paper describes only what is implemented in the StreamSphere /
%% GenABR codebase. Empirical results are deferred to a follow-up paper
%% once paired baseline data has been collected.
%%
%% Compile with:
%%   pdflatex genabr_paper.tex
%%   pdflatex genabr_paper.tex
%%
\documentclass[conference]{IEEEtran}
\IEEEoverridecommandlockouts

% ── Packages ────────────────────────────────────────────────────────────────
\usepackage{cite}
\usepackage{amsmath,amssymb,amsfonts}
\usepackage{algorithmic}
\usepackage{algorithm}
\usepackage{graphicx}
\usepackage{textcomp}
\usepackage{xcolor}
\usepackage{booktabs}
\usepackage{multirow}
\usepackage{url}
\usepackage{hyperref}

\def\BibTeX{{\rm B\kern-.05em{\sc i\kern-.025em b}\kern-.08em
    T\kern-.1667em\lower.7ex\hbox{E}\kern-.125emX}}

\begin{document}

% ── Title ───────────────────────────────────────────────────────────────────
\title{GenABR: A Tiered, Spatially-Aware Adaptive Bitrate Streaming System with
       Generative-Model Escalation for Mobile Video}

\author{
  \IEEEauthorblockN{Bipin Kumar Chaudhary}
  \IEEEauthorblockA{
    Independent Researcher \\
    Bareilly, India \\
    bkumar28899@gmail.com
  }
}

\maketitle

% ── Abstract ────────────────────────────────────────────────────────────────
\begin{abstract}
Conventional Adaptive Bitrate (ABR) controllers respond to throughput
\emph{after} it has already changed. In high-mobility scenarios such as
trains, vehicles, and walking commutes, this reactive design causes visible
quality degradation and rebuffering events when the device crosses
predictable but uncharacterised regions of weak signal. We present
\textbf{GenABR}, a working ABR intelligence layer that operates above an
unmodified HLS.js player and aims to make the decision proactive by
combining (i) a probabilistic spatial look-ahead derived from the user's
current heading and speed, (ii) a continuously updated radio map built from
crowd-sourced and personalised observations, and (iii) a three-tier
inference stack --- on-device Guard, server-side statistical Student, and
LLM-based Oracle --- that escalates only when the cheaper tiers are
uncertain. The system is implemented end-to-end on a public streaming
platform, logs every inference cycle, snaps location data to 200~m tiles
before persistence for privacy, and uses low-power geolocation to bound
battery cost. This paper documents the architecture and the evaluation
framework we have constructed; empirical comparison against reactive
baselines is reserved for a follow-up after paired sessions are collected.
\end{abstract}

\begin{IEEEkeywords}
adaptive bitrate streaming, mobility-aware networking, HLS, tiered
inference, large language models, edge intelligence, quality of experience
\end{IEEEkeywords}

% ============================================================================
\section{Introduction}
% ============================================================================
Mobile video already accounts for the largest share of consumer Internet
traffic, and high-mobility consumption --- on commutes, in vehicles, on
trains --- continues to grow. Standard client-side ABR controllers such as
BOLA and the HLS.js default policy treat available bandwidth as a
quasi-stationary stochastic process estimated from recent segment download
times. This assumption is reasonable in fixed-network or low-mobility
contexts, but it fails systematically when the user crosses a network
``cliff'': a tunnel, a known weak-signal corridor, an inter-cell handover at
speed. By the time the player observes the throughput drop, its buffer has
already begun to drain. Two failure modes follow. The first is a stall ---
the player exhausts its buffer entirely before the next segment is ready.
The second, and more frequent in practice, is an ABR-driven bitrate
downshift, in which the controller cliffs from 720p or 1080p to 360p in
order to avoid the stall, producing a perceptible quality jolt that user
studies consistently rank as a major QoE regression even when no actual
freezing occurs.

The two failures share a single root cause: \emph{the controller has no
information beyond the throughput it has already measured}. We argue that
mobile video QoE for moving users is better modelled as a problem of
\emph{decision under spatial uncertainty} than as one of throughput
forecasting. The spatial distribution of weak-signal regions on a given
road or transit corridor changes slowly, while the throughput history
window over which reactive controllers operate refreshes every few seconds.
A controller that knows where it is heading, how fast, and what the network
has historically looked like along the projected path, can begin to
pre-load video while the connection is still strong and absorb the
upcoming degradation without quality loss.

GenABR pursues this idea with three constraints in mind that distinguish
the work from prior context-aware ABR proposals:
\begin{itemize}
  \item \textbf{Cost discipline.} A predictive controller that consumes
        large-language-model inference on every decision cycle would not
        be deployable. GenABR uses a tiered inference stack so that the
        LLM is invoked only when cheaper layers are uncertain.
  \item \textbf{Privacy.} A system that records full-precision GPS
        traces creates a continuous commute-reconstruction risk. GenABR
        snaps location data to 200~m tile centres at ingestion, matching
        the resolution at which prediction is performed and discarding
        sub-tile precision before it ever lands in storage.
  \item \textbf{Battery.} Continuous high-accuracy GPS is the dominant
        battery cost in typical smartphone web video stacks. GenABR
        sources position from Wi-Fi/cell triangulation, which is roughly
        an order of magnitude lower in power draw and remains adequate
        at the 200~m tile resolution at which the controller operates.
\end{itemize}

This paper documents the system as implemented. It describes the four
architectural pillars, the data model, the three-segment evaluation
harness we have built, and the two specific code-level mitigations for
privacy and battery. It does not claim empirical superiority over reactive
baselines; that comparison requires paired GenABR-on / GenABR-off sessions
on identical routes and is the explicit subject of follow-up work.

% ============================================================================
\section{Background and Motivation}
% ============================================================================

\subsection{HTTP Adaptive Streaming}
HLS \cite{hls} and MPEG-DASH \cite{dash} divide a video into fixed-duration
segments encoded at multiple bitrates. A client-side controller selects the
bitrate of the next segment from a discrete bitrate ladder using observed
throughput, current player buffer occupancy, and a QoE objective. In most
production browser deployments the controller is HLS.js with its default
BOLA-style policy. GenABR is implemented as a buffer-target override above
this player rather than as a replacement controller, so the base bitrate
selection logic is preserved and the system can fall back to standard
behaviour when the override is inactive.

\subsection{The reactive failure mode}
Consider a user walking into a known coverage shadow with a 720p stream and
a 5~s player buffer. As the device crosses the coverage boundary, segment
downloads stretch from $\sim$~1~s each to several seconds. Within a few
seconds the controller observes the slowdown, requests a lower-bitrate
representation, and either (a) succeeds in maintaining playback at a
substantially lower quality for the duration of the shadow, or (b) fails
to fetch even the lower-bitrate segment in time and stalls. In neither
case can the controller exploit the fact that, two minutes earlier, the
device was on the same street with five times the available bandwidth and
could have prefetched the upcoming segments at the higher bitrate.

\subsection{Quality of Experience metrics}
The QoE literature treats three components as primary contributors to
perceived video quality: average video quality (commonly measured by VMAF
\cite{vmaf}), quality variability across the session ($\sigma_{\text{VMAF}}$), and
rebuffering events. Even a single-second stall is often weighted as
equivalent to a multi-point average-quality drop. In the present
implementation, per-frame VMAF is not computed: instead we report a
bitrate-derived VMAF proxy obtained from a fixed lookup table over the
HLS bitrate ladder. We are explicit about this throughout: the metric is
labelled ``VMAF-proxy'' in both the dashboard and the data model, and
replacing it with real Netflix VMAF SDK output is listed as required
future work before publication of empirical claims.

% ============================================================================
\section{Related Work}
% ============================================================================
Reactive ABR has been studied for over a decade. BOLA \cite{bola} casts the
problem as Lyapunov-style optimisation; MPC \cite{mpc} uses receding-horizon
control over a small look-ahead window; Pensieve \cite{pensieve} learns a
controller via deep reinforcement learning. All three derive their state
from recent throughput history alone, and none consult any spatial signal.

Several prior works have introduced location context to ABR. Xia
\textit{et al.} \cite{xia} build per-cell-tower throughput models; Stick
\textit{et al.} \cite{stick} demonstrate that mobility speed is a
significant predictor of throughput variance. These are the closest
points of comparison for our radio-map layer. GenABR extends this line in
three ways: it represents path uncertainty as a probability-weighted cone
rather than a single predicted trajectory, it auto-promotes weak-signal
tiles to dead-zone markers without requiring stalls as triggers, and it
introduces a tiered inference engine in which a generative model serves as
the high-uncertainty escalation layer.

Mobility prediction itself is a mature area, with HMM and deep-learning
trajectory models in widespread use for cellular handover preparation;
GenABR uses a simpler fixed-probability cone, deferring trajectory
learning to future work.

% ============================================================================
\section{System Architecture}
% ============================================================================
GenABR consists of four loosely-coupled components: a frontend telemetry
service running in the browser; a backend inference engine that combines
a probabilistic spatial cone with a shadow radio map; a tiered inference
escalation pipeline; and a structured observation layer that logs every
decision for later analysis. Figure~1 (reserved for a forthcoming
camera-ready revision) summarises the data flow; this section describes
each component as currently implemented in the codebase.

\subsection{Frontend telemetry service}
The frontend is implemented in Angular~19. A telemetry service measures
the device state every \textbf{4~s} and pushes a sample into a local ping
buffer. To avoid storing redundant samples, the service applies
change-detection thresholds: a ping is enqueued only if the buffer level
has changed by at least 2~s, the device has moved at least 50~m, the
bitrate has changed, the network connection type has changed, or 15~s
have elapsed since the last enqueue (a heartbeat). A batch of ten enqueued
pings, or the end-of-session signal, triggers a single HTTP POST to the
backend.

A separate \textbf{prediction service} polls the backend every
\textbf{25~s} for an updated buffer-target recommendation. The first such
request is fired immediately on the first GPS fix to warm the
server-side cache before the steady-state poll cycle begins.

Each ping contains: snapped latitude and longitude (see
Section~\ref{sec:privacy}), speed (km/h), heading, current buffer level
(s), current bitrate (kbps), the most recent downlink and RTT estimates
from the browser's Network Information API, the connection type
(\texttt{4g}, \texttt{3g}, \texttt{2g}, \texttt{slow-2g}, \texttt{wifi}),
and the device battery level when accessible.

\subsection{Probabilistic prediction cone}
\label{sec:cone}
The prediction cone is generated centred on the device's current position
along its current heading. Three branches are evaluated, with fixed
probabilities reflecting that straight-on continuation is most likely and
sharp turns are equally likely on either side:
\begin{itemize}
  \item Straight ($+0\degree$ offset), probability $0.70$
  \item Left ($-30\degree$ offset), probability $0.15$
  \item Right ($+30\degree$ offset), probability $0.15$
\end{itemize}
Each branch is sampled at three look-ahead horizons of $30$, $60$, and
$90$~s, weighted $0.50$, $0.30$, $0.20$ respectively (so points further
in the future contribute proportionally less to the composite risk).
For each sampled point, the system computes the destination tile and
queries the Shadow Map for that tile's fused coverage score
$c(x,y) \in [0, 1]$. The point risk is then
\[
  r_{\text{point}} = p_{\text{branch}} \cdot (1 - c(x, y)) \cdot w_{\text{horizon}}
\]
and the overall spatial risk is the sum of all point risks across all
branches, clipped to $[0, 1]$. We emphasise that the branch probabilities
are presently \emph{fixed constants}, not learned --- replacing them with a
Markov road-topology model is a future-work item.

\subsubsection{Corridor scanner}
In parallel to the cone, a corridor scanner steps along the device's
current heading in 200~m increments up to a maximum of 5~km, querying the
fused score at each step. The first contiguous run of tiles with fused
score below a dead-zone threshold of $0.35$ is identified; the scanner
returns its entry-time-from-now (s), duration (s), required prebuffer
seconds with a $1.30\times$ safety factor, and a feasibility flag set
when the device's current bandwidth is sufficient to download enough
video before reaching the dead zone. The corridor scanner is invoked
only when the device speed exceeds 5~km/h.

\subsection{Shadow Network Map}
\label{sec:shadow}
The Shadow Map is a MongoDB collection (\texttt{radio\_map\_cache}) keyed
on a 200~m geographic tile identifier. Per-tile fields include the
running averages of measured downlink, RTT, and player buffer level; a
sample count; and a fused coverage score. Tiles with no observed samples
fall back to a static crowd-sourced layer (placeholder pluggable to
OpenSignal/OpenCellID); unmapped tiles use a default coverage score of
$0.6$, slightly above the safe threshold and well above the dead-zone
threshold, so that unknown territory does not trigger unnecessary
prebuffering on first visit.

The map is updated by two complementary write paths:
\begin{enumerate}
  \item \textbf{Stall-triggered ingestion.} When the player reports a
        rebuffer event, the (snapped) stall coordinates are written
        directly as a dead-zone marker with a low signal score derived
        from the concurrent downlink reading. This is the conventional
        path used by stall-driven prior work.
  \item \textbf{Auto-promotion at session end.} After every session, the
        \texttt{updateRadioMapFromSession} routine walks all tiles
        touched by the session's pings, applies a Welford-style running
        average to the per-tile downlink statistics, and promotes any
        tile satisfying $\bar{d}_{\text{tile}} < 1.0$~Mbps with sample
        count $\geq 5$ to a dead-zone marker. This path requires no
        stall as a trigger.
\end{enumerate}
The auto-promotion threshold (1.0~Mbps) corresponds to the floor below
which the lowest bitrate ladder rung (360p, $\sim$0.4~Mbps) is at risk;
the sample-count gate (5) prevents single-visit anomalies from
contaminating the map. Stale samples are down-weighted: tiles not visited
within the last 14~days have their effective historical sample count
halved when new data arrives, giving fresh observations more leverage.

\subsection{Tiered inference engine}
\label{sec:tiered}
The escalation hierarchy has three layers, each with progressively higher
latency and cost, and is engineered so that the cheapest tier capable of
handling a given decision actually handles it.

\paragraph{Guard (on-device, free)}
The Guard runs entirely client-side. If the device speed is at or above
3~km/h, the Guard \emph{never} clears the cycle: even on currently
excellent signal, a moving user must be allowed to reach the spatial
look-ahead. If the device is below 3~km/h, the Guard additionally requires
a stable connection type (\texttt{4g} or \texttt{wifi}), a downlink of at
least 5~Mbps, a buffer depth of at least 15~s, and zero stalls in the
preceding 60~s. Only when all four conditions are met does Guard ``pass''
and the cycle is short-circuited with no backend call. The 3~km/h
threshold is set just above the typical apparent-speed noise floor of the
HTML5 Geolocation API for stationary devices ($\approx$~3~km/h with
$\pm$1--3~m accuracy at 1~Hz).

\paragraph{Student (server-side statistical, microcent-level cost)}
When Guard does not pass, the request is dispatched to a Node.js endpoint
\texttt{POST /genabr/decision}. The Student tier computes the spatial
risk from the prediction cone, then applies a \emph{network overlay}
based on the most recent 5 downlink and 5 RTT readings supplied by the
client. The overlay produces a signed risk delta clamped to
$[-0.20, +0.50]$, with structured factor tags drawn from the following
vocabulary: \texttt{conn\_2g}, \texttt{conn\_3g},
\texttt{rtt\_critical} ($>$1000~ms), \texttt{rtt\_high} ($>$500~ms),
\texttt{rtt\_elevated} ($>$200~ms), \texttt{rtt\_rising\_fast},
\texttt{rtt\_rising}, \texttt{dl\_degrading} (with magnitude-scaled
boost), \texttt{dl\_improving}, \texttt{dl\_volatile}, and
\texttt{healthy\_network}. The downlink trend is computed via ordinary
least-squares slope over the last five readings, which correctly
classifies plateau-then-drop patterns that a simple last-minus-first
delta would miss. The composite risk maps to a buffer-target tier:
\begin{itemize}
  \item $\text{risk} < 0.20$: \texttt{normal} (10~s / 30~s caps)
  \item $0.20 \leq \text{risk} < 0.45$: \texttt{prebuffer\_moderate} (20~s / 45~s)
  \item $\text{risk} \geq 0.45$: \texttt{prebuffer\_aggressive} (30~s / 60~s)
\end{itemize}
A Student \emph{confidence} score is computed as the distance to the
nearest decision boundary divided by a half-zone of $0.125$, capped at
$1.0$. When confidence falls below $0.60$ the request is escalated to
the Oracle.

\paragraph{Oracle (cloud LLM, sub-cent cost)}
The Oracle is OpenAI \texttt{gpt-4o-mini} with temperature $0.2$ and a
JSON-output schema. The prompt embeds the spatial cone summary, the
current-tile observed statistics, the recent network readings and trend,
the connection type, and the structured network-factor delta. The model
returns an adjusted risk, a recommendation, the corresponding buffer
caps, a confidence value, and a free-text reasoning string. To bound
latency from the user's perspective, the Oracle is dispatched
asynchronously in the background: the HTTP response returned to the
client is the Student's tentative answer with an
\texttt{oracle\_pending} flag; the next prediction cycle reads the
Oracle's result from a Redis cache and upgrades the recommendation if it
arrived. A per-tile in-flight lock prevents duplicate Oracle calls when
two consecutive cycles both miss the cache during the model's response
window. A per-user limiter caps Oracle calls at ten per minute. When the
Oracle call fails (network error, timeout, schema violation), the system
falls back to a heuristic adjustment computed from speed, downlink
trend, and the time-of-day peak-hours flag, and the failure is logged.

Per-session cost is computed from prompt and completion token counts using
current OpenAI \texttt{gpt-4o-mini} pricing: \$0.15 per million prompt
tokens and \$0.60 per million completion tokens. A baseline cost floor of
\$0.0001 per session is applied at metric computation time so that
zero-Oracle sessions remain numerically comparable to GenABR sessions.

\subsection{Observation layer}
Every Student-tier invocation writes a \texttt{StudentDecision} document,
even when no Oracle is fired, containing: location, speed and speed
category, spatial risk, network delta, structured factor tags, effective
risk, recommendation, confidence, connection type, latest downlink,
latest RTT, and an \texttt{oracle\_pending} boolean. Every Oracle
invocation additionally writes an \texttt{OracleDecision} document
containing the prompt-context tile statistics, the model used, the token
counts, the LLM-returned adjusted risk and reasoning, the heuristic
recommendation that would have applied in the LLM's absence, a
\texttt{diverged} boolean, the corridor-scan results
(\texttt{has\_dead\_zone}, \texttt{dead\_zone\_entry\_sec},
\texttt{dead\_zone\_duration\_sec}, \texttt{corridor\_feasible}), and
the failure-mode flag. \texttt{StudentDecision} records are retained for
30~days; \texttt{OracleDecision} for 90~days.

The \texttt{healthy\_network} tag is logged when no other risk-elevating
factor fires and the device is on \texttt{4g} or \texttt{wifi} with
RTT~$<$~100~ms and downlink~$>$~2~Mbps under a stable trend; recording
this affirmative ``no risk found'' state allows the eventual analysis to
distinguish ``system evaluated and approved'' from ``system was never
invoked''.

% ============================================================================
\section{Privacy Preservation}
\label{sec:privacy}
% ============================================================================
The system records continuous location data, which is a privacy concern
that must be addressed at the implementation layer rather than as a
deployment afterthought. We adopt three measures.

\textbf{Snap-to-tile at ingestion.} The browser geolocation API returns
positions whose precision (under low-power Wi-Fi/cell triangulation) is
typically 50--100~m. The prediction pipeline operates at a 200~m tile
resolution throughout: the corridor scanner steps in 200~m increments,
the prediction cone samples tile-quantised points, the radio-map cache
is keyed on tile identifiers, and dead-zone markers are recorded at
tile-level granularity. Storing the full reported position is therefore
gratuitous data collection. The backend telemetry handler snaps every
incoming latitude/longitude pair to its enclosing tile centre before
writing it to the \texttt{telemetry\_pings} or stall collections; the
GeoJSON \texttt{location} field used by the 2dsphere geospatial index is
populated from the snapped coordinates. No code path after this
snapping step has access to sub-tile precision. Stall events embedded
in \texttt{StreamingSession} (which persist longer than the 2-day
\texttt{TelemetryPing} TTL) are snapped through the same routine.

\textbf{Bounded retention.} \texttt{TelemetryPing} carries a 2-day TTL;
all derived intelligence (radio-map running averages, dead-zone markers,
session summaries) is kept at tile resolution and is not user-keyed.

\textbf{Disabled high-accuracy GPS.} The frontend never requests
\texttt{enableHighAccuracy: true}: the GPS chip is not invoked, and the
data the system never receives at high precision cannot be exfiltrated
or compelled.

These measures are local mitigations rather than a complete privacy
model. Federated radio-map updates with formal differential-privacy
guarantees, and an explicit user opt-in flow, are proposed as follow-up
work.

% ============================================================================
\section{Battery Cost Mitigation}
% ============================================================================
On a typical Android browser, requesting
\texttt{enableHighAccuracy: true} from the Geolocation API engages the
GPS chip directly, with an additional power draw on the order of
50--100~mW when active continuously. Across a 30-minute streaming
session this corresponds to roughly 0.4--0.5\% extra battery on a
typical 4000~mAh device --- bounded but measurable, and entirely
attributable to the predictive layer.

The implementation declines high-accuracy mode and instead obtains
position from the system's Wi-Fi/cell triangulation provider. The
resulting reports are accurate to roughly 50--100~m, which is below the
200~m tile resolution at which all prediction is performed: the cone
sample, corridor step, radio-map key, and dead-zone marker are all
quantised to the same grid, so coarse positioning does not degrade
prediction accuracy. The \texttt{maximumAge} parameter is also raised
from 3~s to 5~s, allowing the system to reuse cached fixes when the
device is briefly stationary and avoiding redundant chip wake-ups.

The remaining contributors --- a 25-second prediction-loop POST and a
batched telemetry POST every 3--5 minutes (after the change-detection
filter described in the previous section drops 60--80\% of redundant
pings) --- are negligible compared to the GPS overhead they replace.

We make no quantitative battery claim in this paper. A controlled
measurement on a reference device, comparing GenABR-on, GenABR-off, and
GenABR-on-with-high-accuracy-GPS, is part of the future work plan.

% ============================================================================
\section{Implementation}
% ============================================================================
The system is implemented as an intelligence layer on top of
\textbf{StreamSphere}, an HLS streaming platform with S3-backed segment
delivery. No modifications are made to the segment delivery path; the
GenABR controller intercepts only the HLS.js buffer-length configuration
keys (\texttt{maxBufferLength}, \texttt{maxMaxBufferLength}), preserving
full CDN performance and enabling clean fallback to the default ABR
policy when the override is inactive. Concretely:

\begin{itemize}
  \item \textbf{Frontend.} Angular~19. Geolocation API, Battery API,
        and Network Information API. HLS.js for the player.
  \item \textbf{Backend.} Node.js with Express, deployed on Vercel.
        MongoDB for persistence (sessions, pings, radio-map cache,
        dead-zone markers, Student/Oracle decision logs). Upstash Redis
        for the Oracle in-flight lock, cached LLM results, the global
        GenABR enable flag, and rate counters.
  \item \textbf{LLM provider.} OpenAI \texttt{gpt-4o-mini} via the
        official Node SDK, called with \texttt{response\_format =
        json\_object}, temperature 0.2, max-tokens 300.
  \item \textbf{Geospatial.} Custom 200~m tile grid; MongoDB
        \texttt{2dsphere} index on the GeoJSON \texttt{location} field
        for dead-zone proximity queries.
\end{itemize}

A \textbf{global GenABR enable flag} is persisted in Redis and exposed
through admin endpoints; toggling it switches all clients to the
unmodified HLS.js policy on their next prediction cycle. This is the
primary control surface for paired A/B data collection.

A \textbf{three-segment admin dashboard} renders, in real time: the
per-segment QoE comparison (Section~\ref{sec:eval}); the Oracle
intelligence panel (decision counts, divergence rate, token totals,
cost, recommendation distribution, network-factor distribution at
trigger time, corridor-scanner statistics including the proactive
lead-time metric); the Student events panel (top-10 network factors,
connection-type distribution, and the share of decisions that escalated
to Oracle); and the recent-sessions log with per-session VMAF-proxy,
$\sigma$, stalls, buffer level, $\Phi$ score, Oracle cost, and route ID.

% ============================================================================
\section{Evaluation Framework}
\label{sec:eval}
% ============================================================================
A direct GenABR-vs-baseline comparison computed over all sessions would
mix three segments with different expected behaviour and produce an
uninterpretable headline number. We therefore classify every session
into one of three non-exclusive segments by aggregating its
\texttt{TelemetryPing} stream:

\begin{itemize}
  \item \textbf{MOBILE} --- $\max(\text{speed}) > 5$~km/h. The device
        was in motion at some point; the spatial look-ahead is the
        relevant mechanism.
  \item \textbf{MOBILE+POOR\_SIGNAL} --- mobile and either average
        downlink $<$~1.5~Mbps or minimum downlink $<$~0.5~Mbps. Both
        the corridor scanner and the network-overlay buffer elevation
        are the relevant mechanisms. This is the highest-value segment
        for the system.
  \item \textbf{STATIONARY+GOOD} --- $\max(\text{speed}) \leq 5$~km/h
        and average downlink $\geq$~3~Mbps. Control group: the Guard
        should pass and the system should defer to HLS.js. Any
        non-zero $\Delta$ in this segment indicates a measurement
        artefact rather than a benefit.
\end{itemize}

Within each segment, sessions are grouped by the
\texttt{genabr\_active} flag and the following per-session metrics are
aggregated:

\begin{itemize}
  \item \textbf{Average VMAF-proxy} --- bitrate-derived quality
        estimate. The current implementation uses a fixed lookup table
        over the bitrate ladder; replacing this with real per-frame
        VMAF SDK output is a prerequisite for any quality claim in a
        camera-ready paper.
  \item \textbf{$\sigma_{\text{VMAF-proxy}}$} --- session-level
        standard deviation. A direct measure of bitrate-switching
        intensity, which is the proximate user-visible quality
        regression we aim to reduce.
  \item \textbf{Stall count and total stall duration} --- the
        rebuffer-event metrics.
  \item \textbf{Average and minimum buffer level} --- the latter is a
        leading indicator of stall risk; the former demonstrates
        proactive buffer fill.
  \item \textbf{Bitrate variance across raw pings} --- complements the
        bitrate-switch event log by capturing micro-fluctuations.
  \item \textbf{Oracle cost (USD)} --- per-session inference cost
        derived from token counts. Reported separately as the cost
        side of the cost-vs-benefit story.
  \item \textbf{Tier distribution} --- counts of Guard / Student /
        Oracle invocations per session, which characterise the
        architectural cost profile.
  \item \textbf{Route identifier} --- a stable hash of the (snapped)
        first and last GPS tile, ordered so that reverse traversal of
        the same physical route yields the same identifier. Allows the
        ``$K$ unique routes'' summary the system produces.
  \item \textbf{Proactive lead time} --- average dead-zone-entry-time
        across cycles where the corridor scanner identified a dead
        zone \emph{and} the inference engine returned a prebuffer
        recommendation. No reactive baseline produces a non-zero value
        here by construction.
\end{itemize}

\subsection{The composite $\Phi$ metric}
We compute a composite QoE-to-cost score
\begin{equation}
  \Phi = \frac{\overline{\text{VMAF-proxy}} - \alpha \, \sigma - \beta \, N_{\text{stall}} \, T_{\text{stall}}}{C_{\text{session}}}
\end{equation}
with $\alpha = 0.5$, $\beta = 20$, and $C_{\text{session}}$ the per-session
inference cost in USD. $C_{\text{session}}$ is floored at \$0.0001 to
keep $\Phi$ finite for zero-Oracle baseline sessions. We surface $\Phi$
on the dashboard as a single diagnostic, but we recommend that the
primary research story be told with the constituent metrics above
plotted on a (cost, QoE) Pareto plane. Composite scores can compress
genuinely informative trade-offs into a single number that obscures
which mechanism is actually contributing.

\subsection{What we will measure}
We do not present empirical results in this paper. The current dataset
contains only GenABR-enabled sessions; without the corresponding
GenABR-disabled sessions on the same routes, no $\Delta$ can be
computed and no claim about the system's behaviour can be supported.
The next phase of work is paired A/B data collection: alternating the
global GenABR enable flag across sessions on a fixed set of routes,
across multiple devices and signal conditions, to populate the three
segment $\times$ two-condition cells with statistically meaningful
sample sizes (target: $n \geq 30$ per cell). Once collected, the
intended primary finding is presented as a (cost, $\Delta$-stalls) and a
(cost, $\Delta\sigma_{\text{VMAF-proxy}}$) Pareto comparison, segmented as
above, alongside a tier-distribution table that characterises
\emph{where} GenABR's cost was incurred.

% ============================================================================
\section{Discussion and Limitations}
% ============================================================================

We list the known limitations of the implementation explicitly so that
the gap between the system as presented and the system as a fully
publishable result is unambiguous.

\textbf{VMAF-proxy, not VMAF.} The current quality metric is a fixed
lookup from bitrate to a VMAF-equivalent. This is acceptable for
internal comparison between configurations of the same player on the
same content, but is not a substitute for per-frame SDK measurement and
should not be reported as ``VMAF'' in any external publication. Real
VMAF requires running the segment files through Netflix's VMAF SDK
(libvmaf via FFmpeg) and persisting per-segment scores; this is a
build-system change rather than a research one.

\textbf{Fixed branch probabilities.} The cone probabilities $0.70 /
0.15 / 0.15$ are hard-coded constants. Replacing them with a Markov
model trained on OpenStreetMap road topology and historical GPS traces
is straightforward in principle and would let the cone narrow on
straight highways and widen at intersections; we have not done it.

\textbf{Cold-start.} The system requires accumulated tile data to
recommend prebuffering for previously unseen regions. A first-time
traversal of a new route falls back to the default coverage of $0.6$
and behaves indistinguishably from baseline ABR. The auto-promotion
threshold of five samples per tile means most regions become
predictable after two or three traversals.

\textbf{Crowd-sourced layer.} The static layer is currently a
placeholder; live ingestion of OpenSignal/OpenCellID grids is a
deployment task, not a research one.

\textbf{Data-plan-aware buffer cap.} The Implementation section's
buffer cap is hard-coded at 300~s. A region-aware cap derived from
local mobile-data-plan minima (e.g., the average daily allowance of
the lower-tier providers in the deployment region) would tighten the
``why not always max-buffer?'' answer; this is a configuration change.

\textbf{Oracle latency tail.} The cloud LLM has a typical response
window of 2--8~s. The shadow-buffer / in-flight-lock design absorbs
this gracefully: the user is never blocked on the LLM. Sustained API
slowness, however, would degrade the 90-second look-ahead guarantee
and is best mitigated by self-hosted inference for the Student tier.

\textbf{Sample size in this paper.} Empirical claims require the
paired-baseline data described in Section~\ref{sec:eval}; we do not
attempt to make any here.

\textbf{Scope.} The system targets web-based HLS playback. Native iOS
and Android applications using AVPlayer / ExoPlayer would require
re-implementation of the buffer-target override mechanism, though the
backend inference engine would be reusable.

% ============================================================================
\section{Future Work}
% ============================================================================

In priority order:

\begin{enumerate}
  \item \textbf{Paired A/B data collection.} Alternate the
        \texttt{genabr\_active} flag across sessions on a fixed route
        set; target $n \geq 30$ per segment per condition.
  \item \textbf{Real VMAF.} Replace the bitrate-lookup proxy with
        per-segment libvmaf SDK output, reported to the database
        alongside the existing fields.
  \item \textbf{Trained branch probabilities.} Replace the fixed
        cone weights with a Markov model trained on OSM topology and
        historical traces.
  \item \textbf{Battery measurement.} Conduct a controlled
        device-side power measurement comparing GenABR-on, GenABR-off,
        and GenABR-on-with-high-accuracy-GPS.
  \item \textbf{Differential-privacy radio-map federation.} Extend
        the snap-to-tile measure with per-update noise and aggregation
        across users to support privacy-preserving cross-user map
        improvement.
  \item \textbf{Static layer ingestion.} Wire the OpenSignal /
        OpenCellID feed into the radio-map cache as the default for
        unmapped tiles.
  \item \textbf{Self-hosted Student model.} Replace the present
        statistical Student with a small distilled model trained on
        the accumulating \texttt{StudentDecision} corpus, removing the
        backend round-trip and lowering tail latency.
  \item \textbf{Native player support.} Port the buffer-target
        override to iOS AVPlayer and Android ExoPlayer.
\end{enumerate}

% ============================================================================
\section{Conclusion}
% ============================================================================

We have described GenABR, an end-to-end implemented ABR intelligence
layer that aims to convert the high-mobility QoE failure mode of
reactive controllers from a measurement problem into a spatial-prediction
problem. The system combines a probabilistic cone, a continuously
updated radio map, and a three-tier inference stack with a generative
model as the high-uncertainty escalation layer; it logs every decision
for analysis; and it incorporates explicit privacy and battery
mitigations at the implementation layer rather than as deployment
afterthoughts. We have not, in this paper, claimed empirical superiority
over any reactive baseline: that comparison requires paired sessions on
identical routes and is the explicit subject of follow-up work. The
contribution presented here is the system, the evaluation harness, and
the three-segment framework within which the eventual comparison can be
made cleanly.

% ── References ──────────────────────────────────────────────────────────────
\begin{thebibliography}{00}
\bibitem{hls} R. Pantos and W. May, ``HTTP Live Streaming,'' IETF RFC 8216, 2017.
\bibitem{dash} ISO/IEC 23009-1, ``Dynamic adaptive streaming over HTTP (DASH),'' 2014.
\bibitem{vmaf} Z. Li, A. Aaron, I. Katsavounidis, A. Moorthy, and M. Manohara, ``Toward a practical perceptual video quality metric,'' Netflix Tech Blog, 2016.
\bibitem{bola} K. Spiteri, R. Urgaonkar, and R. K. Sitaraman, ``BOLA: Near-optimal bitrate adaptation for online videos,'' in \textit{IEEE INFOCOM}, 2016.
\bibitem{mpc} X. Yin, A. Jindal, V. Sekar, and B. Sinopoli, ``A control-theoretic approach for dynamic adaptive video streaming over HTTP,'' in \textit{ACM SIGCOMM}, 2015.
\bibitem{pensieve} H. Mao, R. Netravali, and M. Alizadeh, ``Neural adaptive video streaming with Pensieve,'' in \textit{ACM SIGCOMM}, 2017.
\bibitem{xia} F. Xia \textit{et al.}, ``A measurement study of cellular network performance,'' in \textit{Proc. ACM IMC}, 2014.
\bibitem{stick} J. Stick \textit{et al.}, ``Mobility-aware bandwidth prediction for adaptive video streaming,'' in \textit{Proc. ACM MMSys}, 2016.
\end{thebibliography}

\end{document}
